\section{Introduction}

A realized return does not identify skill. When a decision rule is run on one
price path and turns a profit, the profit confounds at least three sources: the
quality of the decisions, the structural exposure those decisions imply, and the
drift of the path itself. Separating them requires a counterfactual---but not any
counterfactual. The relevant comparison is an implementation that carries the
\emph{same} structural burden on the \emph{same} path, differing only in the margin
whose value is in question. Absent such a comparison, profitability and exposure are
not separately identified, and a backtest, however carefully constructed, remains a
description rather than a test.

This paper develops a general apparatus for that comparison, which we call
\emph{structure-preserving randomization inference}. The idea is a conditional
randomization test in the sense of Fisher \cite{Fisher1935} and Cand\`es et al.\
\cite{Candes2018}: fix the realized structure, and re-randomize only the object of
interest from a known conditional law. Concretely, let a decision sequence be overlaid
on an exogenous process; the realized data induce a \emph{structure}---the geometry of
the sequence---and a placement of that geometry. Holding the structure and the process
fixed while randomizing placement isolates the decisions' marginal contribution, and
its validity rests on the exchangeability of placements under the conditioning law, not
on any model for the process. The construction is indexed by a family of
structure-preserving measures $\mathbb{Q}=\{\mathbb{Q}_\theta:\theta\in\Theta\}$, all
fixing the structure and the path but differing in which additional features of the
realized arrangement they hold fixed; each $\mathbb{Q}_\theta$ defines a sharp null
$H_0^\theta$ and an estimand $q_\theta=\mathbb{Q}_\theta\big(T(S)\ge T(S_0)\big)$, where
$T$ is the statistic and $S_0$ the realized placement. Which features count as
``structure'' and which as ``the decision'' is precisely the analyst's modeling choice,
and making that choice explicit is what turns an informal benchmark into an
identification statement.

The leading application, and the one we develop in full, is the entry-timing skill of a
trading strategy. Here the structure is the realized trade geometry---the price path,
trade count, holding-period multiset, inter-trade-gap multiset, capital weights, and
direction sequence---and the object randomized is the placement of trades on the
calendar. The test asks a narrow, falsifiable question: conditional on its realized
structural burden, does a strategy's entry placement earn more than a structurally
matched random implementation would earn on the same path? A strategy then exhibits
entry-placement skill only when its realized timing beats that counterfactual, not
merely when it is profitable. The empirical testbed is daily gold futures (GC=F) over
2002-03-04 to 2026-04-02, an asset whose buy-and-hold cumulative return of $14.670$ over
the window makes the confound between drift and timing especially stark.

Three features make this a methods paper rather than a study of any one asset. First,
the test is a finite-sample-exact conditional randomization test (Proposition
\ref{prop:q-validity}): under $H_0^\theta$ the realized placement is exchangeable with
its $\mathbb{Q}_\theta$-draws, so the Monte Carlo tail area is super-uniform at every
level, with the standardized conditional excess $\mathrm{RCSI}_z$ retained only as a
descriptive, approximately pivotal effect size. Second, the measure family supports a
partial-identification theory of timing skill (Section~\ref{sec:theory}). Because the
measures differ exactly in where they draw the line between structure and placement, a
rejection under one but not another is confounded with the conditioning choice, which
the data do not pin down; we therefore identify \emph{measure-invariant} entry-placement
skill as rejection of $H_0^\theta$ for every $\theta$ in an admissible class, and we
characterize the test's local power---one-sided Gaussian power
$\beta(h)=\Phi(h-z_{1-\alpha})$ along the Pitman sequence $\zeta_N=h/\sqrt N$, with a
minimum detectable edge that shrinks at the $N^{-1/2}$ rate---and its orthogonality:
against a pure structural-exposure alternative the test carries noncentrality zero, by
construction. Third, that orthogonality is shown head-to-head against return-based
tests. On known-truth worlds (Section~\ref{sec:competitor}), White's Reality Check
\cite{white2000} and Hansen's SPA \cite{hansen2005}---tests of profitability rather than
placement---reject a profitable-but-untimed strategy about two-thirds of the time
(rejection rates $0.660$ and $0.638$), while the structure-preserving null holds its
nominal size ($0.050$). Return-based tests mistake structural exposure for skill; this
test does not.

The paper's \emph{primary contribution} is this framework: the conditional randomization
test, its finite-sample validity, and the partial-identification theory of the measure
family that turns the choice of null into an explicit identification assumption. The
remaining results are \emph{secondary}, present to make the framework usable rather than
to stand alone. A calibration-and-power study---five synthetic worlds with known ground
truth and a positive control on the real GC=F path
(Sections~\ref{sec:simulation-validation} and \ref{sec:positive-control})---shows the
test is correctly sized and that its power rises smoothly with signal quality, so a
non-rejection reflects the strategy rather than a powerless test; the head-to-head
competitor comparison locates precisely where return-based inference fails; and a
cross-asset panel of $322$ strategy-by-asset evaluations across $47$ instruments
(Section~\ref{sec:cross-asset-multiplicity}) applies the test at breadth under joint
multiplicity control. Applied to the GC=F panel, no strategy rejects at $p\le0.05$ before
or after Benjamini--Hochberg adjustment, and the verdict is robust to a leading-gap
feasibility restriction but sensitive to a context-matched measure
(Section~\ref{sec:schedule-sensitivity})---which is exactly the partial-identification
content of the measure family made empirical. The substantive reading is not that the
strategies are bad, but that within the scope of the timing test, realized profitability
is often not separately identifiable from structural exposure.

\section{Related literature}

\paragraph{Data snooping and backtest overfitting.} The literature on strategy
evaluation is largely a literature on multiplicity. White's Reality Check
\cite{white2000}, Hansen's SPA test \cite{hansen2005}, and the universe-search framing of
Sullivan, Timmermann, and White \cite{sullivan1999} control the error of selecting the
best rule \emph{across} a universe of candidates; the backtest-overfitting program
\cite{bailey2017,baileylopez2014,lopezdeprado2018} and the factor-zoo critique of
Harvey, Liu, and Zhu \cite{harveyliuzhu2016} extend the same warning to research degrees
of freedom and the resulting inflation of discoveries. These procedures take each
strategy's realized decision sequence as given and ask whether the \emph{best} of many is
genuinely good. Our object is orthogonal and within-strategy: holding the realized trade
geometry and price path fixed, we ask whether a single rule's entry \emph{placement}
carries information. The two layers are complementary---the between-strategy correction
bounds selection risk, the conditional test isolates timing---and we run the panel under
both.

\paragraph{Resampling of returns.} Bootstrap methods for dependent data---the stationary
bootstrap of Politis and Romano \cite{politisromano1994} and the moving-block bootstrap of
K\"unsch \cite{Kunsch1989}---preserve the statistical properties of the \emph{return
series} and resample it to quantify uncertainty in profitability under a generic
return-generating process. By construction they perturb the path. The structure-preserving
test does the opposite: it holds the realized path fixed and randomizes only placement, so
it inherits the market's autocorrelation, fat tails, volatility clustering, and regime
dynamics without a return model, and answers a timing question that a path-perturbing
resample cannot pose. The contrast is not rhetorical---Section~\ref{sec:competitor} shows
return-based tests reject a profitable-but-untimed strategy where the placement test holds
size.

\paragraph{Conditional and Fisher randomization inference.} The procedure is a conditional
randomization test. Its validity derives, in the tradition of Fisher \cite{Fisher1935}
and the exact and group-invariant theory of Lehmann and Romano \cite{LehmannRomano2005},
from a known re-randomization of the realized data rather than a postulated population
model; Ernst \cite{Ernst2004} surveys this exactness and Pesarin and Salmaso
\cite{PesarinSalmaso2010} treat permutation tests that preserve rich dependence rather
than assuming i.i.d.\ data. The specific posture---condition on a structure, resample only
the object of interest from a known conditional law---is the conditional randomization
construction of Cand\`es et al.\ \cite{Candes2018}. What is new here is the conditioning
set: the realized trade geometry of a strategy, which makes the relevant exchangeability
an exchangeability of \emph{schedules} rather than of \emph{returns}, in the sense of
Aldous \cite{Aldous1985}.

\paragraph{Selective and post-selection inference.} By defining a test conditional on the
structure a strategy actually produced, the method controls a \emph{selective} type-I
error in the sense of Fithian, Sun, and Taylor \cite{FithianSunTaylor2014}: validity is
recovered by characterizing the law of the statistic conditional on the realized selection
event, the same posture as exact post-selection inference \cite{LeeSunSunTaylor2016}. This
deliberately differs from the simultaneous-over-all-selections guarantee of Berk et al.\
\cite{BerkBrownBujaZhangZhao2013}: we condition on the realized structure rather than
protect against every structure that might have been chosen. The measure family makes the
conditioning event itself the modeling object, so that what is held fixed---and therefore
what skill is identified against---is stated explicitly rather than left implicit.

\paragraph{Multiple-testing control.} Applied at breadth across a $322$-test panel,
inference is corrected jointly. We report Benjamini--Hochberg \cite{benjaminihochberg1995}
false-discovery-rate control and the $q$-value perspective of Storey \cite{Storey2002},
the stepwise data-snooping procedure of Romano and Wolf \cite{romanowolf2005}, and the
Benjamini--Yekutieli \cite{BenjaminiYekutieli2001} guarantee that holds under the positive
and, conservatively, arbitrary dependence that correlated assets and overlapping strategy
logic induce.

\paragraph{Efficiency and adaptive markets.} The economic backdrop is the long debate over
whether persistent abnormal performance is extractable from price data. Efficiency
arguments emphasize its difficulty \cite{fama1970}; the adaptive-markets view allows for
episodic, regime-dependent opportunity \cite{lo2004}; and evidence on time-series momentum
documents directional effects in some asset classes while robust inference remains hard
\cite{moskowitz2012}. We take no stance on this debate. The contribution is an
identification tool: a way to ask, for any one strategy on any one path, whether its timing
adds value beyond the exposure it carries.
